% from ega3-1-fr.tex
\oldpage[III]{106}
\begin{proposition}[2.3.5]
\label{III.2.3.5-fr}
Soient $Y$ un préschéma intègre noethérien, $X$ un préschéma intègre,
$f:X\to Y$ un morphisme projectif birationnel. Il existe alors un Idéal
fractionnaire cohérent $\mathcal J\subset\mathcal R(Y)$ (II, 8.1.2) tel
que $X$ soit $Y$-isomorphe au préschéma obtenu en faisant éclater
$\mathcal J$ (II, 8.1.3). En outre, il existe un ouvert $U$ de $Y$ tel que
la restriction de $f$ à $f^{-1}(U)$ soit un isomorphisme de $f^{-1}(U)$
sur $U$ (cf. I, 6.5.5), et que $\mathcal J|U$ soit inversible.
\end{proposition}

\begin{proof}
Comme il existe un $\mathcal O_X$-Module inversible $\mathcal L$ très
ample pour $f$ (II, 4.4.2 et 5.3.2), on peut appliquer
(\hyperref[III.2.3.4.1-fr]{2.3.4.1}), et on voit que $X$ s'identifie à
$\operatorname{Proj}(\mathcal S)$, où
$\mathcal S=\bigoplus_{n\geq0}f_*(\mathcal L^{\otimes n})$. On sait en
outre que les $f_*(\mathcal L^{\otimes n})$ sont des
$\mathcal O_Y$-Modules sans torsion (I, 7.4.5), donc il en est de même du
$\mathcal O_Y$-Module $\mathcal S$, et par suite $\mathcal S$ s'identifie
canoniquement à un sous-$\mathcal O_Y$-Module de
$\mathcal S\otimes_{\mathcal O_Y}\mathcal R(Y)$ (I, 7.4.1)~; ce dernier
est un faisceau \emph{simple} (I, 7.3.6) qui est connu lorsqu'on connaît
sa restriction à un ouvert non vide, par exemple à un ouvert non vide
$U'\subset U$ tel que $\mathcal L|f^{-1}(U')$ soit isomorphe à
$\mathcal O_X|f^{-1}(U')$. Comme par hypothèse les
$f_*(\mathcal L^{\otimes n})|U'$ sont alors isomorphes à
$\mathcal O_Y|U'$, on voit que
$\mathcal S\otimes_{\mathcal O_Y}\mathcal R(Y)$ est un
$\mathcal R(Y)$-Module isomorphe à $\mathcal R(Y)[T]$, où $T$ est une
indéterminée, et $\mathcal S$ est \textbf{(TN)}-isomorphe à la
sous-$\mathcal O_Y$-Algèbre engendrée par l'image canonique de
$f_*(\mathcal L)$ dans $\mathcal S\otimes_{\mathcal O_Y}\mathcal R(Y)$
(\hyperref[III.2.3.4.1-fr]{2.3.4.1})~; mais si on identifie ce dernier
faisceau à $\mathcal R(Y)[T]$, l'image de $f_*(\mathcal L)$ s'identifie à
$\mathcal J T$, où $\mathcal J$ est un sous-$\mathcal O_Y$-Module
cohérent (\hyperref[III.2.2.1-fr]{2.2.1}) de $\mathcal R(Y)$, dont la
restriction à $U'$ est isomorphe à $\mathcal O_Y|U'$, et qui par suite
est tel que $\mathcal J|U$ soit inversible. On voit alors que $\mathcal S$
est \textbf{(TN)}-isomorphe à $\bigoplus_{n\geq0}\mathcal J^n$, ce qui
achève la démonstration.
\end{proof}

\begin{corollary}[2.3.6]
\label{III.2.3.6-fr}
Sous les hypothèses de (\hyperref[III.2.3.5-fr]{2.3.5}), supposons en
outre que, pour tout sous-$\mathcal O_Y$-Module cohérent
$\mathcal J\neq0$ de $\mathcal R(Y)$, il existe un
$\mathcal O_Y$-Module inversible $\mathcal L$ tel que
\[
\Gamma(Y,\mathcal L\otimes_{\mathcal O_Y}
\mathcal{H}om(\mathcal J,\mathcal O_Y))\neq0~;
\]
alors, dans l'énoncé de (\hyperref[III.2.3.5-fr]{2.3.5}), on peut supposer
que $\mathcal J$ est un Idéal de $\mathcal O_Y$. Cette condition
supplémentaire est toujours vérifiée s'il existe un
$\mathcal O_Y$-Module ample.
\end{corollary}

\begin{proof}
En effet, on a (0\textsubscript{I}, 5.4.2)
\[
\mathcal L\otimes\mathcal{H}om(\mathcal J,\mathcal O_Y)
=\mathcal{H}om(\mathcal L^{-1},\mathcal{H}om(\mathcal J,\mathcal O_Y))
=\mathcal{H}om(\mathcal J\otimes\mathcal L^{-1},\mathcal O_Y)~;
\]
l'hypothèse signifie donc qu'il y a un homomorphisme non nul $u$ de
$\mathcal J\otimes\mathcal L^{-1}$ dans $\mathcal O_Y$. Comme, pour tout
$y\in Y$, $(\mathcal J\otimes\mathcal L^{-1})_y$ s'identifie à un
sous-$\mathcal O_y$-Module du corps des fractions $\mathcal R(Y)_y$ de
$\mathcal O_y$ (I, 7.1.5), $u_y$ est nécessairement injectif, donc $u$
est un isomorphisme de $\mathcal J\otimes\mathcal L^{-1}$ sur un Idéal
$\mathcal J'$ de $\mathcal O_Y$. Mais comme
\[
\operatorname{Proj}\left(\bigoplus_{n\geq0}\mathcal J^n\right)
\quad\text{et}\quad
\operatorname{Proj}\left(
\bigoplus_{n\geq0}(\mathcal J\otimes\mathcal L^{-1})^n\right)
\]
sont $Y$-isomorphes (II, 3.1.8), cela prouve la première assertion du
corollaire. Pour démontrer la seconde, notons que
$\mathcal F=\mathcal{H}om_{\mathcal O_Y}(\mathcal J,\mathcal O_Y)$ est
cohérent et $\neq0$, puisqu'il existe un ouvert $U$ de $Y$ tel que
$\mathcal J|U$ soit inversible. Si $\mathcal L$ est un
$\mathcal O_Y$-Module ample, il existe un entier $n$ tel que
$\mathcal F(n)=\mathcal F\otimes\mathcal L^{\otimes n}$ soit engendré par
ses sections au-dessus de $Y$ (II, 4.5.5)~; \emph{a fortiori}, on a
$\Gamma(Y,\mathcal F(n))\neq0$, d'où la conclusion.
\end{proof}

\begin{corollary}[2.3.7]
\label{III.2.3.7-fr}
Soient $X$ et $Y$ deux schémas intègres, projectifs sur un corps $k$, et
soit $f:X\to Y$ un $k$-morphisme birationnel. Alors $X$ est $k$-isomorphe
à un $Y$-schéma obtenu en faisant éclater un sous-schéma fermé $Y'$
(non nécessairement réduit) de $Y$.
\end{corollary}

